\section{Commitment bases and compound commitments}
\label{app:bases}

Section~\ref{sec:functional-calm} considers only the output events of
Definition~\ref{def:output-event}, each of which applies at most the one seal
its output requires.  Real systems declare which commitments
(Definition~\ref{def:commitment-basis}) they may make, may apply several in one
output event, may commit more than an output requires, and may restrict when
each commitment is available.  This appendix shows that none of this changes
the results: compound commitments behave like single ones, and every system's
runs are contained in the runs of Section~\ref{sec:functional-calm}, which we
call canonical.

\begin{definition}[Commitment basis]
\label{def:basis}
A \emph{commitment basis} $\mathcal B$ is a set of commitments available to a
system.
\end{definition}

Commitments are \emph{partial} functions, so a system may authorize a
commitment only at states satisfying a precondition of its own.  In
Figure~\ref{fig:election-frames}(d), the poll-closing commitment is defined at
every state; a different system might authorize it only after some candidate
has a majority.  A system whose polls cannot close has no such commitment in its
basis.

\begin{definition}[Compound output event]
\label{def:compound-output-event}
Let $(i,R)$ be a state and $o\in R(i)$.  An output event $\Out{\gamma}{o}$
under a basis $\mathcal B$ applies a finite, possibly empty sequence
$\gamma=c_1\cdots c_n$ of commitments from $\mathcal B$, each $c_m$ defined at
the state its predecessors leave, and then emits $o$, all atomically;
$\varepsilon$ is the empty sequence.  If $\gamma$ leaves remainder table $R'$,
the event is \emph{legal} at $(i,R)$ when every remaining row's output
\outrel $o$:
\begin{equation}
  (k,u)\in R' \Longrightarrow o\OutputOrd u,
  \quad\text{or equivalently}\quad R'\subseteq R[\mathrm{out}\OutputGe o].
  \label{eq:supported-output-action}
\end{equation}
\end{definition}

Every commitment maps a state to a state and leaves the current input
unchanged, so $R'$ still contains the current row $(i,o)$.

\begin{definition}[Canonical basis, canonical event]
\label{def:canonical-commitment-basis}
The \emph{canonical commitment basis} consists of the seals of
Definition~\ref{def:output-event},
\begin{equation}
  \mathcal B_S
    =\{s_o:o\in O,\ \operatorname{dom}(s_o)\neq\varnothing\},
  \label{eq:canonical-commitment-basis}
\end{equation}
where $s_o$ is defined at the states $(i,R)$ with $(i,o)\in R$ and
$R[\mathrm{out}\OutputGe o]\neq R$.  At $(i,R)$, let $\kappa_o$ be $s_o$ if
$s_o$ is defined there and $\varepsilon$ otherwise.  The \emph{canonical event}
for $o$ is $\Out{\kappa_o}{o}$: the output event $\mathsf{output}(o)$ of
Definition~\ref{def:output-event}.
\end{definition}

Three facts relate an arbitrary output event to the canonical one.

\begin{proposition}[Output events]
\label{prop:canonical-output-preservation}
Let $(i,R)$ be a state and $o\in R(i)$.
\begin{enumerate}
  \item \emph{Containment.}  Under any commitment basis, a legal output event
        $\Out{\gamma}{o}$ takes $(i,R)$ to some $(i,R')$ with
        $R'\subseteq R[\mathrm{out}\OutputGe o]$.
  \item \emph{Canonical support.}  The canonical event $\Out{\kappa_o}{o}$ is
        legal at $(i,R)$ and takes it to $(i,R[\mathrm{out}\OutputGe o])$.  It
        therefore retains every row that any legal $\Out{\gamma}{o}$ retains.
  \item \emph{Commitment-free events are canonical.}  If $\Out{\varepsilon}{o}$
        is legal, it takes $(i,R)$ to $(i,R[\mathrm{out}\OutputGe o])$; it is
        legal exactly when $o$ is exposable.
\end{enumerate}
\end{proposition}

\begin{proof}
Containment is Equation~\ref{eq:supported-output-action}, and canonical
support is Equation~\ref{eq:output-action}.  The event $\Out{\varepsilon}{o}$
leaves $R$ unchanged, so it is legal exactly when
$R\subseteq R[\mathrm{out}\OutputGe o]$, that is, when
$R=R[\mathrm{out}\OutputGe o]$, which is Definition~\ref{def:exposable}.
\end{proof}

\begin{lemma}[Compound commitments are commitments]
\label{lem:compound-commitment}
The composite of a nonempty sequence of commitments, defined wherever the
sequence is, is a commitment.
\end{lemma}

\begin{proof}
Each step leaves the current input unchanged and maps a state to a state, so the
composite does too, and it removes a nonempty set of rows because the first
step does.  Suppose step $m$ removes $(k,u)$ from table $R_m$.  No row of $R_m$
at $k$ \outrel $u$, and the final table is contained in $R_m$, so no remaining
row at $k$ \outrel $u$: Equation~\ref{eq:commitment-excludes} holds.
\end{proof}

A compound output event therefore has the same effect as an event that
applies the composite commitment alone, and
Proposition~\ref{prop:canonical-output-preservation} covers it as stated.  In
a functional specification the composite of seals is again a seal.

\begin{definition}[Run under a basis]
\label{def:run-under-basis}
Fix a commitment basis $\mathcal B$.  A \emph{run under $\mathcal B$} is a
finite transition sequence from a canonical state $\sigma_S(i)$ whose output
events $\Out{\gamma}{o}$ use only commitments from $\mathcal B$.  The runs of
Definition~\ref{def:run}, the \emph{canonical runs}, are the runs under
$\mathcal B_S$ whose output events are all canonical.
\end{definition}

Our results quantify over canonical runs rather than over the runs of any
particular system.  To see that this excludes no system's behavior, call two
runs \emph{equivalent} when they start at the same canonical state and agree
once commitment sequences are erased: wherever one has $\mathsf{input}(j)$ the
other does, and wherever one has $\Out{\gamma}{o}$ the other has
$\Out{\gamma'}{o}$ for some $\gamma'$.  They accept the same inputs and emit
the same outputs, in the same order, which is all the environment observes.
Equivalent runs may differ in the commitments their output events apply, and
hence in their remainder tables.

\begin{lemma}[Canonical equivalents]
\label{lem:canonical-equivalent}
Every run under any commitment basis is equivalent to a canonical run.  After
each event, the canonical run's remainder table contains the original run's.
\end{lemma}

\begin{proof}
By induction on the events.  Both runs start at the same canonical state, so
their remainder tables are initially equal.  Suppose that before some event the
original run's table $R$ and the canonical run's table $\widehat R$ satisfy
$R\subseteq\widehat R$; we show that the inclusion holds afterward.  An input
event $\mathsf{input}(j)$ that is legal at $R$ is also legal at $\widehat R$.
It retires every row whose input does not \inrelbase $j$, keeping
$R[\mathrm{in}\InputGe j]$ in the original run and
$\widehat R[\mathrm{in}\InputGe j]$ in the canonical run.  Since
$R\subseteq\widehat R$, the first remains contained in the second.  For a legal
output event $\Out{\gamma}{o}$ with $o\in R(i)$ taking $R$ to $R'$, legality
(Equation~\ref{eq:supported-output-action}) gives
$R'\subseteq R[\mathrm{out}\OutputGe o]$.  Since
$o\in R(i)\subseteq\widehat R(i)$, canonical support
(Proposition~\ref{prop:canonical-output-preservation}) applies at $\widehat R$:
$\Out{\kappa_o}{o}$, with $\kappa_o$ taken at $(i,\widehat R)$, takes
$\widehat R$ to $\widehat R[\mathrm{out}\OutputGe o]$, which contains
$R[\mathrm{out}\OutputGe o]$.
\end{proof}

The converse does not hold: a canonical run need not be equivalent to any run
under a particular system's basis.  That basis may lack the commitment an
output event requires, and so disable the output altogether.  Restricting
attention to one system's runs could therefore hide non-monotonicity.  This is
why Consistency (Definition~\ref{def:consistency}) is defined over canonical
runs rather than over a system's own.

\paragraph{Application-specific commitments.}
Locks and timestamps are commitments of this kind: two-phase locking commits an
order among some transactions, and timestamp assignment places one transaction
relative to all others.  Section~\ref{sec:ambiguity-story} illustrates both,
and Appendix~\ref{app:cc-bases} writes the two bases down.  By
Lemma~\ref{lem:canonical-equivalent}, the runs of either are contained in
canonical runs.
